%% Generated by Sphinx.
\def\sphinxdocclass{report}
\documentclass[a4paper,11pt,english]{sphinxmanual}
\ifdefined\pdfpxdimen
   \let\sphinxpxdimen\pdfpxdimen\else\newdimen\sphinxpxdimen
\fi \sphinxpxdimen=.75bp\relax
\ifdefined\pdfimageresolution
    \pdfimageresolution= \numexpr \dimexpr1in\relax/\sphinxpxdimen\relax
\fi
\newdimen\sphinxremdimen\sphinxremdimen = 11pt
%% let collapsible pdf bookmarks panel have high depth per default
\PassOptionsToPackage{bookmarksdepth=5}{hyperref}
%% turn off hyperref patch of \index as sphinx.xdy xindy module takes care of
%% suitable \hyperpage mark-up, working around hyperref-xindy incompatibility
\PassOptionsToPackage{hyperindex=false}{hyperref}
%% memoir class requires extra handling
\makeatletter\@ifclassloaded{memoir}
{\ifdefined\memhyperindexfalse\memhyperindexfalse\fi}{}\makeatother

\PassOptionsToPackage{booktabs}{sphinx}
\PassOptionsToPackage{colorrows}{sphinx}

\PassOptionsToPackage{warn}{textcomp}

\catcode`^^^^00a0\active\protected\def^^^^00a0{\leavevmode\nobreak\ }
\usepackage{cmap}
\usepackage{fontspec}
\defaultfontfeatures[\rmfamily,\sffamily,\ttfamily]{}
\usepackage{amsmath,amssymb,amstext}
\usepackage{polyglossia}
\setmainlanguage{english}



\setmainfont{FreeSerif}[
  Extension      = .otf,
  UprightFont    = *,
  ItalicFont     = *Italic,
  BoldFont       = *Bold,
  BoldItalicFont = *BoldItalic
]
\setsansfont{FreeSans}[
  Extension      = .otf,
  UprightFont    = *,
  ItalicFont     = *Oblique,
  BoldFont       = *Bold,
  BoldItalicFont = *BoldOblique,
]
\setmonofont{FreeMono}[Scale=0.9,
  Extension      = .otf,
  UprightFont    = *,
  ItalicFont     = *Oblique,
  BoldFont       = *Bold,
  BoldItalicFont = *BoldOblique,
]



\usepackage[Bjarne]{fncychap}
\usepackage{sphinx}

\fvset{fontsize=auto}
\usepackage{geometry}


% Include hyperref last.
\usepackage{hyperref}
% Fix anchor placement for figures with captions.
\usepackage{hypcap}% it must be loaded after hyperref.
% Set up styles of URL: it should be placed after hyperref.
\urlstyle{same}

\addto\captionsenglish{\renewcommand{\contentsname}{Contents}}

\usepackage{sphinxmessages}
\setcounter{tocdepth}{1}

\usepackage{microtype}

\title{Jobs — Market Analysis}
\date{Oct 03, 2026}
\release{1.0}
\author{Benoit Chevallier-Mames}
\newcommand{\sphinxlogo}{\vbox{}}
\renewcommand{\releasename}{Release}
\makeindex
\begin{document}

\pagestyle{empty}
\sphinxmaketitle
\pagestyle{plain}
\sphinxtableofcontents
\pagestyle{normal}
\phantomsection\label{\detokenize{index::doc}}


\sphinxstepscope


\chapter{Market Analysis — Self\sphinxhyphen{}Hosted Job\sphinxhyphen{}Search Aggregators}
\label{\detokenize{analysis:market-analysis-self-hosted-job-search-aggregators}}\label{\detokenize{analysis::doc}}
\sphinxAtStartPar
\sphinxstylestrong{Author:} compiled for Benoit Chevallier\sphinxhyphen{}Mames
\sphinxstylestrong{Date:} October 2026
\sphinxstylestrong{Repo under review:} \sphinxcode{\sphinxupquote{sandboxed\sphinxhyphen{}repos/jobs}} (147 companies, Playwright + LLM scoring)


\section{Executive summary}
\label{\detokenize{analysis:executive-summary}}
\sphinxAtStartPar
The self\sphinxhyphen{}hosted job\sphinxhyphen{}aggregator space is \sphinxstylestrong{crowded but fragmented}. There are
dozens of open\sphinxhyphen{}source scrapers on GitHub, most of them maintained by a single
author for their own job hunt. None of them combine all four of the following
properties, which is what makes the repo under review distinctive:


\begin{savenotes}\sphinxattablestart
\sphinxthistablewithglobalstyle
\centering
\begin{tabulary}{\linewidth}[t]{TT}
\sphinxtoprule
\begin{varwidth}[t]{\sphinxcolwidth{1}{2}}
\sphinxstyletheadfamily \sphinxAtStartPar
Property
\sphinxbeforeendvarwidth
\end{varwidth}%
&\begin{varwidth}[t]{\sphinxcolwidth{1}{2}}
\sphinxstyletheadfamily \sphinxAtStartPar
Common in competitors?
\sphinxbeforeendvarwidth
\end{varwidth}%
\\
\sphinxmidrule
\sphinxtableatstartofbodyhook\begin{varwidth}[t]{\sphinxcolwidth{1}{2}}
\sphinxAtStartPar
100+ ATS sources (not just LinkedIn)
\sphinxbeforeendvarwidth
\end{varwidth}%
&\begin{varwidth}[t]{\sphinxcolwidth{1}{2}}
\sphinxAtStartPar
Rare
\sphinxbeforeendvarwidth
\end{varwidth}%
\\
\sphinxhline\begin{varwidth}[t]{\sphinxcolwidth{1}{2}}
\sphinxAtStartPar
Per\sphinxhyphen{}user state (reject/like/apply/K)
\sphinxbeforeendvarwidth
\end{varwidth}%
&\begin{varwidth}[t]{\sphinxcolwidth{1}{2}}
\sphinxAtStartPar
Common
\sphinxbeforeendvarwidth
\end{varwidth}%
\\
\sphinxhline\begin{varwidth}[t]{\sphinxcolwidth{1}{2}}
\sphinxAtStartPar
LLM\sphinxhyphen{}based fit scoring with a profile
\sphinxbeforeendvarwidth
\end{varwidth}%
&\begin{varwidth}[t]{\sphinxcolwidth{1}{2}}
\sphinxAtStartPar
Rare
\sphinxbeforeendvarwidth
\end{varwidth}%
\\
\sphinxhline\begin{varwidth}[t]{\sphinxcolwidth{1}{2}}
\sphinxAtStartPar
Senior\sphinxhyphen{}level focus (SENIORITY\_XP, IC\sphinxhyphen{}level badges)
\sphinxbeforeendvarwidth
\end{varwidth}%
&\begin{varwidth}[t]{\sphinxcolwidth{1}{2}}
\sphinxAtStartPar
Essentially unique
\sphinxbeforeendvarwidth
\end{varwidth}%
\\
\sphinxbottomrule
\end{tabulary}
\sphinxtableafterendhook\par
\sphinxattableend\end{savenotes}

\sphinxAtStartPar
The closest named competitors are \sphinxstylestrong{JobSpy} (consumer boards),
\sphinxstylestrong{ats\sphinxhyphen{}scrapers / jobhive} (ATS\sphinxhyphen{}native, 50+ platforms), \sphinxstylestrong{JobSentinel}
(local\sphinxhyphen{}first UX, preference\sphinxhyphen{}based scoring), and \sphinxstylestrong{JobAggregator}
(FastAPI dashboard tuned for India + remote AI). Each one dominates a
different corner of the design space — none covers the combination this
repo targets.


\section{Methodology}
\label{\detokenize{analysis:methodology}}
\sphinxAtStartPar
We searched GitHub + the open web (October 2026) for projects tagged
\sphinxcode{\sphinxupquote{job\sphinxhyphen{}scraper}}, \sphinxcode{\sphinxupquote{job\sphinxhyphen{}aggregator}}, \sphinxcode{\sphinxupquote{job\sphinxhyphen{}search}}, and \sphinxcode{\sphinxupquote{self\sphinxhyphen{}hosted}}, plus
Scrapfly’s August 2026 roundup of actively\sphinxhyphen{}maintained open\sphinxhyphen{}source
scrapers. For each candidate we extracted:
\begin{itemize}
\item {} 
\sphinxAtStartPar
Source coverage (which boards / ATS)

\item {} 
\sphinxAtStartPar
Deployment model (CLI, FastAPI, Docker, GitHub Pages)

\item {} 
\sphinxAtStartPar
State tracking (does the user have persistent rejected/liked state?)

\item {} 
\sphinxAtStartPar
Scoring (preference rules, LLM, nothing)

\item {} 
\sphinxAtStartPar
Target user (generic job seeker, senior engineer, specific vertical)

\item {} 
\sphinxAtStartPar
Last commit (as a proxy for whether the project still works)

\end{itemize}

\sphinxAtStartPar
Projects with no commits in the past 6 months were excluded unless they
scored exceptionally high on another axis.


\section{Direct competitors}
\label{\detokenize{analysis:direct-competitors}}

\subsection{JobSpy (speedyapply/JobSpy)}
\label{\detokenize{analysis:jobspy-speedyapply-jobspy}}\begin{itemize}
\item {} 
\sphinxAtStartPar
\sphinxstylestrong{Stars:} 4.4k · \sphinxstylestrong{Language:} Python · \sphinxstylestrong{Last active:} Feb 2026

\item {} 
\sphinxAtStartPar
\sphinxstylestrong{Coverage:} LinkedIn, Indeed, Glassdoor, Google Jobs, ZipRecruiter,
Bayt, Naukri, BDJobs

\item {} 
\sphinxAtStartPar
\sphinxstylestrong{Deployment:} \sphinxcode{\sphinxupquote{pip install python\sphinxhyphen{}jobspy}}, returns a pandas DataFrame

\item {} 
\sphinxAtStartPar
\sphinxstylestrong{State tracking:} none — it’s a library, not an app

\item {} 
\sphinxAtStartPar
\sphinxstylestrong{Scoring:} none

\item {} 
\sphinxAtStartPar
\sphinxstylestrong{Target:} developers building their own aggregator

\end{itemize}

\sphinxAtStartPar
\sphinxstylestrong{Verdict:} The de\sphinxhyphen{}facto library for \sphinxstyleemphasis{consumer boards}. Where this repo
focuses on ATS\sphinxhyphen{}native boards (where the hiring companies actually post
first), JobSpy focuses on the aggregators (where the same job appears
late and sometimes reposted). The two are complementary, not competitors.


\subsection{ats\sphinxhyphen{}scrapers / jobhive}
\label{\detokenize{analysis:ats-scrapers-jobhive}}\begin{itemize}
\item {} 
\sphinxAtStartPar
\sphinxstylestrong{Coverage:} 50+ ATS platforms (Greenhouse, Lever, Ashby, Workday, …)

\item {} 
\sphinxAtStartPar
\sphinxstylestrong{Deployment:} library, no API key required, 27 columns of output

\item {} 
\sphinxAtStartPar
\sphinxstylestrong{State tracking:} none

\item {} 
\sphinxAtStartPar
\sphinxstylestrong{Scoring:} none

\item {} 
\sphinxAtStartPar
\sphinxstylestrong{Target:} data engineers wanting raw ATS feeds

\end{itemize}

\sphinxAtStartPar
\sphinxstylestrong{Verdict:} The closest technical sibling. It covers the same universe
of ATS platforms this repo does, but exposes them as a pipeline — no UI,
no per\sphinxhyphen{}user state. If this repo had to “stand on the shoulders of giants”
to go generic, ats\sphinxhyphen{}scrapers would be the dependency to adopt.


\subsection{JobSentinel (cboyd0319/JobSentinel)}
\label{\detokenize{analysis:jobsentinel-cboyd0319-jobsentinel}}\begin{itemize}
\item {} 
\sphinxAtStartPar
\sphinxstylestrong{Coverage:} multiple boards + ATS, with ghost\sphinxhyphen{}job detection

\item {} 
\sphinxAtStartPar
\sphinxstylestrong{Deployment:} local desktop workspace, Python + TypeScript

\item {} 
\sphinxAtStartPar
\sphinxstylestrong{State tracking:} yes — save, notes, contacts, reminders, offers

\item {} 
\sphinxAtStartPar
\sphinxstylestrong{Scoring:} preference\sphinxhyphen{}based (salary floor, keyword match) — the
v1.3 release docs explicitly flag “Extract skills from description”
as deferred pending AI integration

\item {} 
\sphinxAtStartPar
\sphinxstylestrong{Target:} job seeker who wants a private, offline\sphinxhyphen{}first workflow

\end{itemize}

\sphinxAtStartPar
\sphinxstylestrong{Verdict:} The product closest in \sphinxstyleemphasis{spirit} to this repo. Shares the
“private by default, local only, no telemetry” axis. Where this repo
has invested is \sphinxstylestrong{LLM scoring already working} (both Ollama and
Claude paths) and \sphinxstylestrong{big\sphinxhyphen{}tech\sphinxhyphen{}specific seniority badges} (L5/L6, E5/E6,
IC7, …). JobSentinel compensates with a richer UX around applications,
reminders, and ghost\sphinxhyphen{}job scoring — this repo has none of that.


\subsection{JobAggregator (SammyUrfen/JobAggregator)}
\label{\detokenize{analysis:jobaggregator-sammyurfen-jobaggregator}}\begin{itemize}
\item {} 
\sphinxAtStartPar
\sphinxstylestrong{Deployment:} FastAPI dashboard

\item {} 
\sphinxAtStartPar
\sphinxstylestrong{Coverage:} job\sphinxhyphen{}board scrapers + remote\sphinxhyphen{}work APIs + company ATS

\item {} 
\sphinxAtStartPar
\sphinxstylestrong{State tracking:} dedupe + expiration (jobs that disappear from a
board are marked expired)

\item {} 
\sphinxAtStartPar
\sphinxstylestrong{Scoring:} none visible

\item {} 
\sphinxAtStartPar
\sphinxstylestrong{Target:} remote\sphinxhyphen{}first / India\sphinxhyphen{}based systems + AI roles

\end{itemize}

\sphinxAtStartPar
\sphinxstylestrong{Verdict:} Shows a reasonable packaging path: FastAPI + themed
dashboard. The dedupe\sphinxhyphen{}across\sphinxhyphen{}sources + expire\sphinxhyphen{}when\sphinxhyphen{}removed logic is
analogous to this repo’s \sphinxcode{\sphinxupquote{job\_index.json}} + orphan rendering.


\subsection{Find\sphinxhyphen{}Me\sphinxhyphen{}Job (MohamedMamdouh18/Find\sphinxhyphen{}Me\sphinxhyphen{}Job)}
\label{\detokenize{analysis:find-me-job-mohamedmamdouh18-find-me-job}}\begin{itemize}
\item {} 
\sphinxAtStartPar
\sphinxstylestrong{Coverage:} LinkedIn, RemoteOK (daily scrape)

\item {} 
\sphinxAtStartPar
\sphinxstylestrong{Deployment:} self\sphinxhyphen{}hosted pipeline → Notion database + Telegram

\item {} 
\sphinxAtStartPar
\sphinxstylestrong{State tracking:} yes, via Notion

\item {} 
\sphinxAtStartPar
\sphinxstylestrong{Scoring:} \sphinxstylestrong{LLM\sphinxhyphen{}based (0\sphinxhyphen{}100) against CV}, auto\sphinxhyphen{}generates
tailored cover letters for high\sphinxhyphen{}scoring matches

\item {} 
\sphinxAtStartPar
\sphinxstylestrong{Target:} individuals who want end\sphinxhyphen{}to\sphinxhyphen{}end automation up to the
“we drafted the cover letter for you” step

\end{itemize}

\sphinxAtStartPar
\sphinxstylestrong{Verdict:} The only project found that uses an LLM for fit scoring,
like this repo does. The novelty is the auto\sphinxhyphen{}cover\sphinxhyphen{}letter step — a
reasonable next feature for this repo if the user wants to go that way.


\subsection{JobSpy\sphinxhyphen{}API (rainmanjam/jobspy\sphinxhyphen{}api)}
\label{\detokenize{analysis:jobspy-api-rainmanjam-jobspy-api}}\begin{itemize}
\item {} 
\sphinxAtStartPar
\sphinxstylestrong{Stars:} 379 · \sphinxstylestrong{Deployment:} Dockerized, HTTP API, API\sphinxhyphen{}key auth,
rate limiting, proxy support

\item {} 
\sphinxAtStartPar
\sphinxstylestrong{Coverage:} inherits JobSpy’s boards

\item {} 
\sphinxAtStartPar
\sphinxstylestrong{Target:} teams wanting to call a shared backend

\end{itemize}

\sphinxAtStartPar
\sphinxstylestrong{Verdict:} The reference “how to productize an open\sphinxhyphen{}source scraper”
project. The HTTP + proxies + rate limits pattern is what this repo
would need before hosting a SaaS variant.


\subsection{job\sphinxhyphen{}board\sphinxhyphen{}aggregator (Feashliaa/)}
\label{\detokenize{analysis:job-board-aggregator-feashliaa}}\begin{itemize}
\item {} 
\sphinxAtStartPar
\sphinxstylestrong{Scale:} 1,000,000+ positions · 20,000+ companies

\item {} 
\sphinxAtStartPar
\sphinxstylestrong{Platforms:} Greenhouse, Lever, Ashby, Workday, iCIMS, BambooHR,
Paylocity

\item {} 
\sphinxAtStartPar
\sphinxstylestrong{Deployment:} GitHub Actions ETL + client\sphinxhyphen{}side filterable search

\item {} 
\sphinxAtStartPar
\sphinxstylestrong{Workers tuned per platform:} 50 for Workday, 30 for GH/Lever/iCIMS,
10 for BambooHR, 5 for Ashby/Paylocity

\end{itemize}

\sphinxAtStartPar
\sphinxstylestrong{Verdict:} The scale this repo could reach if it stopped curating 147
boards and started indexing every board on each ATS. The tradeoff:
coverage vs. signal. This repo’s 147 are hand\sphinxhyphen{}picked for the user’s
industries (cybersecurity, big\sphinxhyphen{}tech, music tech); Feashliaa’s 20k
companies include fast\sphinxhyphen{}food chains and gyms.


\section{Commercial job\sphinxhyphen{}search CRMs (adjacent category)}
\label{\detokenize{analysis:commercial-job-search-crms-adjacent-category}}
\sphinxAtStartPar
Not direct competitors — different buyer — but worth knowing about:
\begin{itemize}
\item {} 
\sphinxAtStartPar
\sphinxstylestrong{Huntr} — purpose\sphinxhyphen{}built job\sphinxhyphen{}search CRM, Chrome extension to one\sphinxhyphen{}click
save postings, AI resume + cover letter writer

\item {} 
\sphinxAtStartPar
\sphinxstylestrong{Teal} — bookmarked → applied → interviewing → offer stages,
Chrome extension pulling from 40+ boards

\item {} 
\sphinxAtStartPar
\sphinxstylestrong{Careerflow.ai} — LinkedIn profile audit + CRM + outbound tracking

\item {} 
\sphinxAtStartPar
\sphinxstylestrong{Careerkit} — custom status columns, attach the resume version
you submitted per entry

\item {} 
\sphinxAtStartPar
\sphinxstylestrong{JibberJobber} — the “veteran” option, broad job\sphinxhyphen{}search management

\end{itemize}

\sphinxAtStartPar
All of these are web SaaS, require signup, and are optimized for the
\sphinxstyleemphasis{application tracking} phase. This repo’s distinction is that it
operates one phase earlier: \sphinxstylestrong{surfacing postings worth considering}
across hundreds of career pages at once.


\section{Where this repo sits on the market map}
\label{\detokenize{analysis:where-this-repo-sits-on-the-market-map}}
\begin{sphinxVerbatim}[commandchars=\\\{\}]
                            STATE (per\PYGZhy{}user tracking)
                                      ▲
                                      │
                 Teal, Huntr,         │       THIS REPO
                 Careerflow,          │       (147 boards,
                 Careerkit            │        LLM scoring,
                 (SaaS, consumer)     │        local\PYGZhy{}first)
                                      │
                                      │       JobSentinel
                                      │       Find\PYGZhy{}Me\PYGZhy{}Job
                                      │       JobAggregator
                                      │
   ◀──────────────────────────────────┼──────────────────────────────────▶
   NO AGGREGATION                     │                      HEAVY AGGREGATION
                                      │
                                      │       JobSpy, ats\PYGZhy{}scrapers,
                                      │       Feashliaa aggregator,
                                      │       jobhive, Levergreen
                                      │
                                      ▼
                            PIPELINE (no state, no UI)
\end{sphinxVerbatim}

\sphinxAtStartPar
The upper\sphinxhyphen{}right quadrant — \sphinxstylestrong{aggregation + state + LLM scoring +
senior\sphinxhyphen{}level bias} — is sparsely populated. This repo lives there.
Its closest neighbours are JobSentinel (same UX bias, less aggregation)
and Find\sphinxhyphen{}Me\sphinxhyphen{}Job (same LLM\sphinxhyphen{}scoring bias, far less aggregation).


\section{Design decisions that are rare in the field}
\label{\detokenize{analysis:design-decisions-that-are-rare-in-the-field}}
\sphinxAtStartPar
Based on this survey, the following choices in this repo are either
unique or very rare among open\sphinxhyphen{}source competitors:
\begin{enumerate}
\sphinxsetlistlabels{\arabic}{enumi}{enumii}{}{.}%
\item {} 
\sphinxAtStartPar
\sphinxstylestrong{Per\sphinxhyphen{}company dedicated fetchers alongside the generic ATS fetchers.}
Most projects use the generic Greenhouse/Ashby/Workday APIs. This
repo has dedicated code paths for Cisco (\sphinxcode{\sphinxupquote{fetch\_cisco}}), Bose
(\sphinxcode{\sphinxupquote{fetch\_bose}}), IBM (\sphinxcode{\sphinxupquote{fetch\_ibm}}), LinkedIn (public guest pages),
Apple, Google, Meta, Microsoft, NVIDIA — because those sources either
don’t expose a clean API or have schema quirks the generic fetcher
misses. This is maintenance cost but it’s also the reason the repo
has signal where competitors have blanks.

\item {} 
\sphinxAtStartPar
\sphinxstylestrong{Big\sphinxhyphen{}tech\sphinxhyphen{}specific seniority mapping.} \sphinxcode{\sphinxupquote{SENIORITY\_XP{[}Google{]}{[}L7{]} → "\textasciitilde{}13y (Senior Staff)"}}, \sphinxcode{\sphinxupquote{IC\_LEVEL\_XP{[}Netflix{]}{[}L6{]} → "\textasciitilde{}9y (Staff)"}}.
No other open\sphinxhyphen{}source project surveyed does this.

\item {} 
\sphinxAtStartPar
\sphinxstylestrong{Spontaneous\sphinxhyphen{}application surfacing.} The ✉ Spontaneous link on
every company section lets the user send a cold application when the
current board has nothing matching. No competitor does this.

\item {} 
\sphinxAtStartPar
\sphinxstylestrong{Score + reason + persistent cache.} Scores from the LLM get
saved to \sphinxcode{\sphinxupquote{claude\_fit\_cache.json}} with a timestamp, and the paste
bar updates them in place without losing older scores.

\item {} 
\sphinxAtStartPar
\sphinxstylestrong{In\sphinxhyphen{}browser R/⚙/C controls that POST to a local HTTP server.}
Classical SPAs talk to a hosted backend. This repo’s UI talks to a
\sphinxcode{\sphinxupquote{127.0.0.1:8765}} HTTP server co\sphinxhyphen{}located with the Python process.
Simpler, no auth, no CORS pain — only works for a single user.

\end{enumerate}


\section{Realistic paths forward}
\label{\detokenize{analysis:realistic-paths-forward}}

\subsection{Path A — Open\sphinxhyphen{}source as\sphinxhyphen{}is, no hosting}
\label{\detokenize{analysis:path-a-open-source-as-is-no-hosting}}\begin{itemize}
\item {} 
\sphinxAtStartPar
Publish on GitHub with a self\sphinxhyphen{}hosted “quick start” README

\item {} 
\sphinxAtStartPar
Target: senior engineers who already run a dev environment

\item {} 
\sphinxAtStartPar
Install burden: \sphinxcode{\sphinxupquote{git clone + pip install playwright + python3 jobs.py}}

\item {} 
\sphinxAtStartPar
Cost: zero

\item {} 
\sphinxAtStartPar
Distribution: GitHub topic tagging (\sphinxcode{\sphinxupquote{self\sphinxhyphen{}hosted}}, \sphinxcode{\sphinxupquote{job\sphinxhyphen{}search}},
\sphinxcode{\sphinxupquote{job\sphinxhyphen{}aggregator}}) + mention on \sphinxcode{\sphinxupquote{r/cscareerquestions}} + HN “Show HN”

\end{itemize}

\sphinxAtStartPar
\sphinxstylestrong{Difficulty:} Low. Mostly a refactor pass (which the user has already
asked for: separate engine from data, example config, etc.).


\subsection{Path B — Dockerized self\sphinxhyphen{}hosted}
\label{\detokenize{analysis:path-b-dockerized-self-hosted}}\begin{itemize}
\item {} 
\sphinxAtStartPar
Package Playwright + Python + the engine as a \sphinxcode{\sphinxupquote{docker compose up}}

\item {} 
\sphinxAtStartPar
Target: homelabbers (same audience as Immich / Paperless\sphinxhyphen{}ngx)

\item {} 
\sphinxAtStartPar
Install burden: edit \sphinxcode{\sphinxupquote{user\_config.py}}, mount \sphinxcode{\sphinxupquote{./data}}, \sphinxcode{\sphinxupquote{docker up}}

\item {} 
\sphinxAtStartPar
Cost: zero to the user

\item {} 
\sphinxAtStartPar
Distribution: AwesomeSelfHosted list, r/selfhosted, LinuxServer

\end{itemize}

\sphinxAtStartPar
\sphinxstylestrong{Difficulty:} Medium. Playwright headless Chromium in Docker is a
known\sphinxhyphen{}painful area (apt packages, font issues, image size \textasciitilde{}1.5 GB).
The engine already runs headless so no architectural change — only
packaging.


\subsection{Path C — SaaS}
\label{\detokenize{analysis:path-c-saas}}\begin{itemize}
\item {} 
\sphinxAtStartPar
Multi\sphinxhyphen{}tenant, cloud\sphinxhyphen{}hosted, pay\sphinxhyphen{}per\sphinxhyphen{}month

\item {} 
\sphinxAtStartPar
Target: engineers who don’t want to self\sphinxhyphen{}host anything

\end{itemize}

\sphinxAtStartPar
\sphinxstylestrong{Blockers:}
\begin{itemize}
\item {} 
\sphinxAtStartPar
\sphinxstylestrong{Legal.} LinkedIn / Google / Apple scrape protections. hiQ v.
LinkedIn (2022) narrowed the “public data” shield; a hosted service
that scrapes LinkedIn on behalf of third parties is clearly riskier
than running the same code on your own laptop with your own IP.

\item {} 
\sphinxAtStartPar
\sphinxstylestrong{Anti\sphinxhyphen{}bot.} Datacenter IPs get flagged by LinkedIn within hours.
Residential proxies cost \$5\sphinxhyphen{}15/GB and would blow the margin on any
reasonable subscription price.

\item {} 
\sphinxAtStartPar
\sphinxstylestrong{Playwright at scale.} 147 boards × Chromium per fetch × N users is
expensive compute. A full run on the user’s Mac takes 200s with
parallelism and warm cache.

\item {} 
\sphinxAtStartPar
\sphinxstylestrong{State isolation.} JSON state files become multi\sphinxhyphen{}tenant SQL. UI
grows auth, billing, account management.

\end{itemize}

\sphinxAtStartPar
\sphinxstylestrong{Difficulty:} High. The engineering is doable but the regulatory +
cost + anti\sphinxhyphen{}bot story is a sustained headwind.


\subsection{Path D — “Managed self\sphinxhyphen{}hosted”}
\label{\detokenize{analysis:path-d-managed-self-hosted}}\begin{itemize}
\item {} 
\sphinxAtStartPar
The hosted service helps the user DEPLOY the self\sphinxhyphen{}hosted version
(one\sphinxhyphen{}click deploy to Fly.io / Railway / Render)

\item {} 
\sphinxAtStartPar
The user owns the IP, the data, and the risk

\item {} 
\sphinxAtStartPar
The service collects a one\sphinxhyphen{}time setup fee + optional support

\end{itemize}

\sphinxAtStartPar
\sphinxstylestrong{Difficulty:} Low once Path B is done. The pattern is
\sphinxcode{\sphinxupquote{railway.app/template/<yourtemplate>}} with the user’s env vars.


\section{Recommendation}
\label{\detokenize{analysis:recommendation}}
\sphinxAtStartPar
\sphinxstylestrong{Go Path A → Path B.} Open\sphinxhyphen{}source first, Docker second. SaaS is a
legal + operational minefield that extracts value this project
doesn’t need to capture.

\sphinxAtStartPar
Specifically:
\begin{enumerate}
\sphinxsetlistlabels{\arabic}{enumi}{enumii}{}{.}%
\item {} 
\sphinxAtStartPar
Finish the refactor the user already started
(engine/data separation, delete PROFILE.md, \sphinxcode{\sphinxupquote{user\_config.example.py}}).

\item {} 
\sphinxAtStartPar
Write a 15\sphinxhyphen{}line README showing a copy\sphinxhyphen{}paste quick\sphinxhyphen{}start.

\item {} 
\sphinxAtStartPar
Push to GitHub, tag with topics, Show HN once.

\item {} 
\sphinxAtStartPar
Add a \sphinxcode{\sphinxupquote{Dockerfile}} + \sphinxcode{\sphinxupquote{docker\sphinxhyphen{}compose.yml}} in the following month.

\item {} 
\sphinxAtStartPar
Add a one\sphinxhyphen{}click deploy template to Fly.io for users who want
remote\sphinxhyphen{}but\sphinxhyphen{}still\sphinxhyphen{}self\sphinxhyphen{}hosted.

\item {} 
\sphinxAtStartPar
Skip SaaS entirely. Point users who ask for one to Huntr / Teal
for the application\sphinxhyphen{}tracking side and keep this repo as the
aggregation + scoring side.

\end{enumerate}


\section{Open questions worth asking the user}
\label{\detokenize{analysis:open-questions-worth-asking-the-user}}\begin{itemize}
\item {} 
\sphinxAtStartPar
Target audience: senior engineers only, or do you want it to work
equally well for a mid\sphinxhyphen{}level or a product manager?

\item {} 
\sphinxAtStartPar
SENIORITY\_XP opinionatedness: these are tuned for “\textasciitilde{}5\sphinxhyphen{}15 years into
tech”. Someone who’s 25 years in won’t recognize themselves in the
grid.

\item {} 
\sphinxAtStartPar
Industry breadth: 147 companies currently skew cybersecurity /
big\sphinxhyphen{}tech / music\sphinxhyphen{}tech. For a generic release, a \sphinxcode{\sphinxupquote{\sphinxhyphen{}\sphinxhyphen{}industry cybersecurity}} CLI flag + bundled industry configs would reduce
onboarding friction.

\item {} 
\sphinxAtStartPar
Hosting ambition: are you optimizing for your own hunt only, or
genuinely wanting to help others?

\end{itemize}



\renewcommand{\indexname}{Index}
\printindex
\end{document}